\phantomsection
\chapter*{Résumé étendu en français}\adjustmtc
\addcontentsline{toc}{chapter}{Résumé étendu en français}
\markboth{Résumé étendu en français}{Résumé étendu en français}

\section*{Contexte général et motivations}

Les véhicules autonomes et connectés occupent aujourd'hui une place centrale dans l'évolution des systèmes de transport intelligents. Leur promesse est double : améliorer la sécurité routière en réduisant la dépendance aux décisions humaines, et augmenter l'efficacité du trafic grâce à la coopération entre véhicules, infrastructures et systèmes embarqués. Cette évolution repose cependant sur une condition essentielle : le véhicule doit disposer, à chaque instant, d'une information fiable sur son propre état dynamique et sur l'état de son environnement proche. Les grandeurs telles que la position, la vitesse, l'accélération, la distance inter-véhicules, les forces de contact pneu-chaussée ou les états transmis par les véhicules voisins ne sont pas toujours directement mesurables avec la précision, la fréquence et la robustesse nécessaires. Elles doivent donc être reconstruites à partir de capteurs imparfaits, de modèles physiques et de données échangées par communication.

Dans ce contexte, l'estimation d'état constitue un élément fondamental de la chaîne de décision et de contrôle. Un observateur bien conçu permet de reconstruire des variables non mesurées, de filtrer le bruit, de détecter des incohérences et d'alimenter les lois de commande avec des informations plus régulières que les mesures brutes. Cette fonction devient critique dans les architectures de conduite automatisée, où les décisions de freinage, d'accélération ou de maintien de trajectoire dépendent directement de la qualité des estimations. Une erreur d'estimation peut se propager jusqu'au contrôle longitudinal ou latéral, provoquer une mauvaise évaluation de la distance de sécurité ou masquer l'apparition d'un défaut. La problématique n'est donc pas seulement académique : elle est directement liée à la sûreté de fonctionnement et à l'acceptabilité industrielle des véhicules autonomes.

Les véhicules connectés ajoutent une difficulté supplémentaire. Les fonctions coopératives telles que l'Adaptive Cruise Control (ACC), le Cooperative Adaptive Cruise Control (CACC) ou le peloton de véhicules utilisent les informations reçues par communication véhicule-à-véhicule (V2V). Ces informations peuvent enrichir la perception locale, réduire les temps de réaction et améliorer la fluidité du trafic. En revanche, elles introduisent aussi une surface d'attaque et de défaillance. Les messages peuvent être retardés, perdus, corrompus ou falsifiés. Une attaque par injection de fausses données peut, par exemple, modifier l'accélération transmise par un véhicule leader et perturber la réponse de ses suiveurs. De même, une interruption de communication ou un défaut capteur peut rendre une donnée partagée inutilisable. Les mécanismes d'estimation doivent donc être capables de distinguer les informations fiables des informations douteuses, de reconstruire les perturbations inconnues et de limiter leur effet sur la commande.

Cette thèse s'inscrit dans cette problématique générale de résilience. Le terme résilience désigne ici la capacité d'un système d'estimation et de contrôle à maintenir un comportement acceptable malgré les incertitudes de modèle, les défauts, les attaques et les limitations de communication. La thèse ne se limite pas à détecter qu'une anomalie est présente. Elle cherche à reconstruire son effet, à réduire son influence dans la boucle de contrôle et à fournir une estimation exploitable pour les fonctions de conduite. Cette orientation explique le choix d'un cadre fondé sur les observateurs à entrées inconnues, les observateurs hybrides intégrant l'apprentissage, et les mécanismes distribués de confiance.

\section*{Objectifs scientifiques de la thèse}

L'objectif principal de la thèse est de développer de nouveaux algorithmes d'estimation avancée pour améliorer la résilience des véhicules autonomes et connectés. Trois axes scientifiques structurent le manuscrit. Le premier concerne la reconstruction d'attaques ou de perturbations modélisées comme des entrées inconnues dans un système CACC. Le deuxième traite des incertitudes non linéaires et variables dans le temps à l'aide d'une architecture hybride, combinant un observateur fondé sur le modèle et un observateur neuronal adaptatif. Le troisième étend l'estimation au cadre distribué des pelotons de véhicules, où les informations reçues des voisins doivent être pondérées selon leur niveau de confiance.

Le premier objectif consiste à concevoir un observateur capable d'estimer simultanément l'état d'un véhicule et les signaux malveillants affectant les communications. Les observateurs à entrées inconnues sont particulièrement adaptés à cette tâche, car ils permettent de traiter une perturbation non mesurée comme une variable à reconstruire. Toutefois, les conditions classiques d'existence de ces observateurs, notamment les conditions de rang et de couplage entre entrées inconnues et sorties mesurées, ne sont pas toujours satisfaites par les modèles réalistes de véhicules. La thèse cherche donc à proposer des formulations discrètes et des transformations adaptées pour rendre cette reconstruction possible dans des cas pratiques.

Le deuxième objectif est lié aux limites des modèles physiques. Un modèle de véhicule utilisé pour la conception d'un observateur doit rester suffisamment simple pour être exploitable en temps réel, mais suffisamment riche pour représenter les phénomènes dominants. Or les forces pneu-chaussée, les variations d'adhérence, les effets aérodynamiques ou les dynamiques d'actionneurs peuvent être difficiles à décrire avec précision. Les méthodes purement apprenantes peuvent approximer ces effets, mais elles sont souvent difficiles à certifier et à intégrer dans une architecture sûre. La thèse adopte donc une position intermédiaire : utiliser l'apprentissage pour approximer les dynamiques résiduelles, tout en conservant une structure d'observateur qui fournit des garanties de stabilité et de bornitude.

Le troisième objectif porte sur la coopération. Dans un peloton, chaque véhicule peut utiliser ses mesures locales et les messages reçus de ses voisins. Cette redondance peut améliorer l'observabilité et la robustesse, mais elle devient dangereuse si certains voisins transmettent des données compromises. La thèse propose alors une architecture d'observateur distribué couplée à un mécanisme de confiance. Les poids attribués aux informations reçues ne sont pas fixés une fois pour toutes : ils évoluent en fonction de critères de cohérence portant sur la vitesse, la distance, l'accélération et les estimations globales. L'objectif est de réduire automatiquement l'influence des données non fiables sans rompre la coopération entre véhicules sains.

\section*{Cadre méthodologique}

La méthodologie générale repose sur la redondance analytique. Au lieu de dépendre uniquement de capteurs physiques coûteux ou de messages V2V supposés parfaits, la thèse exploite les relations dynamiques entre les variables du système. Les modèles de véhicule, les équations d'espacement, les sorties mesurées et les lois de commande fournissent des contraintes qui permettent de vérifier la cohérence des données et de reconstruire les grandeurs manquantes. Cette redondance analytique est utilisée comme un capteur logiciel : elle produit des estimations d'état et d'attaques à partir des mesures disponibles.

Le manuscrit combine plusieurs niveaux de modélisation. Le niveau local décrit la dynamique d'un véhicule, ses états longitudinaux ou latéraux, ses entrées de commande et ses sorties mesurées. Le niveau coopératif décrit les interactions entre véhicules, notamment les distances relatives, les vitesses échangées et les informations transmises par communication. Le niveau cyber-physique introduit des attaques ou des défauts sous forme d'entrées inconnues, de délais, de pertes de paquets ou de données incohérentes. Ce choix permet d'étudier la résilience non seulement comme une propriété de robustesse au bruit, mais comme une réponse structurée à des dégradations physiques et informationnelles.

Sur le plan algorithmique, la thèse utilise des observateurs à entrées inconnues, des conditions de synthèse sous forme d'inégalités matricielles linéaires, des transformations de systèmes et des mécanismes d'adaptation neuronale. Les inégalités matricielles linéaires permettent de formuler la stabilité et l'atténuation des perturbations sous une forme calculable numériquement. Les transformations de modèles permettent de contourner certaines restrictions classiques de rang. Les réseaux neuronaux sont utilisés non pas comme un remplacement complet du modèle physique, mais comme un module d'approximation des résidus. Enfin, les scores de confiance servent à adapter les poids d'un observateur distribué en fonction de la qualité estimée des informations reçues.

La validation est menée à plusieurs niveaux. Des simulations MATLAB/Simulink permettent d'évaluer précisément les erreurs d'estimation, la reconstruction d'attaques et la réponse des observateurs dans des scénarios contrôlés. Des environnements plus proches de l'expérimentation, comme CARLA et Quanser QLabs, introduisent des contraintes supplémentaires : dynamique plus réaliste, bruit, trajectoires complexes, communication simulée et architecture logicielle modulaire. Le dernier chapitre prépare le passage vers l'embarqué en décrivant une plateforme matérielle dédiée à l'acquisition, à la communication et à l'exécution future des algorithmes proposés.

\section*{Estimation résiliente et reconstruction d'attaques pour le CACC}

La première contribution majeure concerne l'estimation résiliente dans un système CACC soumis à des cyberattaques. Dans une architecture CACC, un véhicule suiveur utilise non seulement ses propres mesures, mais aussi les informations envoyées par le véhicule leader ou par les véhicules voisins. Cette coopération permet d'améliorer la réponse longitudinale, de réduire les distances de sécurité et de limiter les oscillations dans le peloton. Toutefois, si une information transmise est falsifiée, le contrôleur peut agir sur une base erronée. Une fausse accélération communiquée peut ainsi provoquer un freinage ou une accélération inadaptée, dégrader la stabilité de chaîne et créer un risque de collision.

La thèse modélise ces signaux corrompus comme des entrées inconnues affectant le modèle discret du système. L'enjeu est de reconstruire ces entrées en même temps que les états du véhicule. Une approche de simple détection fournirait un indicateur d'anomalie, mais ne donnerait pas nécessairement la valeur du signal à compenser. La reconstruction permet au contraire d'estimer l'amplitude et l'évolution temporelle de l'attaque. Cette information peut ensuite être utilisée dans une stratégie de défense, par exemple en soustrayant l'effet estimé de l'attaque ou en adaptant la donnée transmise avant son utilisation par le contrôleur.

Le chapitre développe une structure d'observateur à entrées inconnues en temps discret. La formulation tient compte du fait que les données de contrôle et de communication sont échantillonnées. Cette caractéristique est importante pour l'implémentation automobile, car les calculateurs embarqués fonctionnent avec des périodes d'échantillonnage fixes. Une attention particulière est portée à la reconstruction avec délai borné : l'estimation de l'attaque ne doit pas arriver trop tard, car une compensation tardive pourrait être inutile pour la sécurité. La thèse propose également une approche de discrétisation exacte afin de limiter les erreurs dues au passage du modèle continu au modèle discret.

Les résultats de simulation montrent que l'observateur proposé est capable de reconstruire des attaques variables dans le temps et de maintenir une estimation fiable des états longitudinaux. Les scénarios étudiés incluent des attaques par injection de fausses données et des perturbations de communication. Les comparaisons avec des approches plus classiques mettent en évidence l'intérêt de la formulation proposée, en particulier lorsque la dynamique de l'attaque évolue rapidement. La stratégie de défense améliore la capacité du véhicule à suivre la trajectoire ou l'espacement désiré malgré la présence de données compromises.

Cette contribution établit le socle cyber-résilient de la thèse. Elle montre que l'attaque ne doit pas seulement être considérée comme une anomalie externe, mais comme une variable dynamique à estimer. Elle fournit aussi une première illustration de la redondance analytique : les relations entre les états, les entrées et les sorties permettent de reconstruire une information qui n'est pas directement mesurée.

\section*{Observateurs hybrides et apprentissage de dynamiques non modélisées}

La deuxième contribution traite le cas des dynamiques non linéaires et mal connues. Les véhicules réels présentent des comportements difficiles à modéliser exactement. Les forces latérales, les forces longitudinales, l'adhérence, les effets de charge ou les variations de route introduisent des termes non linéaires qui peuvent changer au cours du temps. Un observateur conçu sur un modèle simplifié peut donc produire une estimation biaisée si ces effets ne sont pas correctement pris en compte.

Pour répondre à cette difficulté, la thèse développe une architecture d'apprentissage fondée sur un observateur à entrées inconnues généralisé et un observateur neuronal adaptatif. L'idée principale est de séparer deux fonctions complémentaires. Le premier niveau, fondé sur le modèle, fournit une estimation robuste et bornée des états ou des résidus. Le second niveau apprend les composantes non modélisées à partir de ces informations. Cette stratégie évite de confier toute l'estimation à un réseau neuronal isolé. Le réseau apprend dans un cadre structuré, supervisé par un observateur qui exploite les équations physiques du système.

Le chapitre commence par analyser les limitations des observateurs à entrées inconnues classiques lorsque la condition de rang n'est pas satisfaite. Une formulation utilisant les dérivées de sortie et une régularisation est ensuite introduite pour construire un modèle généralisé. Cette transformation rend possible l'estimation dans des situations où les méthodes standard seraient trop restrictives. La synthèse de l'observateur est formulée au moyen de conditions matricielles assurant la stabilité de l'erreur d'estimation et un niveau de robustesse vis-à-vis des perturbations.

Le module neuronal est ensuite intégré pour approximer les dynamiques résiduelles. Les paramètres d'apprentissage sont ajustés à partir d'une fonction de perte liée aux erreurs ou aux résidus observés. Une discrétisation de l'observateur neuronal et un calcul de gradients par analyse de sensibilité sont proposés pour rendre la méthode exploitable numériquement. La stratégie de learning continu vise à mettre à jour l'approximation au fil du fonctionnement, tout en évitant une adaptation non contrôlée qui pourrait dégrader la stabilité.

L'application à un modèle de véhicule met en évidence l'intérêt de cette architecture hybride. L'observateur généralisé traite les difficultés structurelles liées aux entrées inconnues, tandis que le module neuronal améliore l'estimation des termes complexes, comme les résidus associés aux forces pneu-chaussée. Les simulations montrent une amélioration de la qualité d'estimation par rapport à une approche purement modèle lorsque les dynamiques non modélisées deviennent significatives. Elles confirment également que l'apprentissage peut être introduit dans un cadre compatible avec les exigences de stabilité et de robustesse.

Cette contribution répond à un enjeu important pour l'industrialisation des méthodes d'estimation. Les modèles analytiques sont interprétables et certifiables, mais parfois incomplets. Les modèles appris sont flexibles, mais difficiles à garantir. L'approche proposée cherche à combiner les avantages des deux familles : l'observateur fournit la structure et la stabilité, tandis que l'apprentissage apporte l'adaptabilité nécessaire face aux incertitudes complexes.

\section*{Estimation distribuée et gestion de la confiance dans les pelotons}

La troisième contribution porte sur les pelotons de véhicules connectés. Dans un peloton, chaque véhicule peut améliorer son estimation en utilisant les informations issues de ses voisins. Cette coopération est utile lorsque certaines grandeurs ne sont pas observables localement ou lorsque les mesures locales sont bruitées. Néanmoins, elle crée une dépendance vis-à-vis de la qualité des messages reçus. Si un voisin transmet une information compromise, son influence peut se propager dans l'observateur distribué et contaminer l'estimation de plusieurs véhicules.

La thèse propose une architecture d'observateur distribué dans laquelle les poids attribués aux voisins sont modulés par un score de confiance. Ce score n'est pas seulement une valeur abstraite. Il est construit à partir d'indicateurs de cohérence entre les données reçues, les mesures locales, les estimations globales et les contraintes physiques du peloton. Les critères incluent notamment la cohérence de vitesse, la cohérence de distance et la cohérence d'accélération. Lorsqu'une donnée reçue devient incompatible avec les informations disponibles localement ou avec la dynamique attendue, son niveau de confiance diminue.

Le mécanisme de confiance est conçu à plusieurs niveaux. La confiance locale évalue la cohérence immédiate des messages reçus. La confiance distribuée tient compte des estimations partagées et de la cohérence avec l'état global reconstruit. Une mémoire temporelle permet de ne pas réagir uniquement à un écart instantané, qui pourrait provenir d'un bruit ou d'un transitoire normal, mais d'accumuler l'historique de fiabilité. Cette mémoire est importante pour distinguer une perturbation ponctuelle d'un comportement durablement suspect. Les scores obtenus sont ensuite normalisés et transformés en poids utilisés par l'observateur.

La conception des poids vise à préserver deux propriétés contradictoires. D'une part, il faut réduire l'influence des véhicules non fiables pour éviter la propagation des attaques. D'autre part, il faut conserver suffisamment de connectivité pour que l'observateur distribué reste utile. Une coupure trop brutale de toutes les communications pourrait dégrader l'estimation, surtout si certaines grandeurs dépendent de l'information coopérative. Le compromis proposé consiste à adapter progressivement les poids en fonction de la confiance, avec des bornes et des règles de normalisation.

Les simulations de peloton montrent que cette approche améliore la résilience face à différents scénarios d'attaque. Lorsqu'un véhicule envoie des données incohérentes, son influence diminue dans les estimations des autres véhicules. Le peloton conserve ainsi une meilleure cohérence d'état et une meilleure capacité de contrôle. Les résultats soulignent également que la confiance doit être liée à des critères physiques et non uniquement à des statistiques de communication. Un message peut être reçu correctement sur le plan réseau tout en étant physiquement faux ; l'observateur doit donc vérifier sa compatibilité avec le comportement dynamique attendu.

Cette contribution élargit la thèse du véhicule isolé vers le système multi-agent. Elle montre que la résilience ne dépend pas seulement de l'observateur local, mais aussi de la manière dont les informations sont partagées, évaluées et pondérées. Elle ouvre ainsi la voie à des architectures de coopération plus sûres pour les véhicules connectés.

\section*{Validation expérimentale simulée et préparation à l'embarqué}

La dernière partie du manuscrit s'intéresse au passage des algorithmes vers des environnements plus proches de l'expérimentation. Les validations purement numériques sont indispensables pour analyser les erreurs et comparer les méthodes, mais elles ne suffisent pas à évaluer les contraintes d'intégration. Un système embarqué doit gérer des fréquences d'échantillonnage, des délais, des pertes de communication, des capteurs hétérogènes et une architecture logicielle modulaire. C'est pourquoi la thèse introduit des plateformes de simulation et de prototypage telles que Quanser QCar2/QLabs et une carte électronique embarquée dédiée.

L'environnement QLabs permet de mettre en place des scénarios de conduite avec véhicules virtuels, capteurs simulés, communication V2V et injection d'attaques. Il sert de banc de validation intermédiaire entre MATLAB/Simulink et les essais réels. Les fonctions de perception, de communication, de contrôle longitudinal et de contrôle latéral peuvent y être testées dans un cadre plus intégré. Cette étape est importante pour vérifier que les observateurs développés ne restent pas isolés, mais peuvent s'insérer dans une chaîne complète comprenant acquisition, estimation, décision et commande.

Le chapitre décrit également une architecture de contrôle combinant le suivi longitudinal, le maintien de trajectoire et des mécanismes coopératifs. L'objectif est de préparer l'utilisation des estimations dans une boucle de conduite réelle. Les données produites par les observateurs et les indicateurs de confiance doivent pouvoir être transmises aux modules de commande, mais aussi enregistrées pour analyse. Cette logique de journalisation facilite la comparaison entre simulation, essais virtuels et futurs essais physiques.

La carte électronique embarquée proposée constitue une base matérielle pour l'implémentation future des observateurs. Elle est pensée comme un noeud d'estimation et de communication capable d'acquérir des données locales, de recevoir des messages V2V, de synchroniser les informations et d'exécuter des traitements de confiance ou d'estimation. La conception prend en compte les contraintes d'un environnement embarqué : alimentation, interfaces capteurs, communication sans fil, routage de circuit imprimé, compatibilité électromagnétique et modularité logicielle. Même si l'intégration complète des algorithmes sur cette carte reste une perspective, sa conception montre que la thèse anticipe le transfert vers des plateformes réelles.

Cette partie expérimentale et matérielle complète les contributions théoriques. Elle met en évidence les contraintes qui devront être traitées pour passer d'un observateur validé en simulation à un module embarqué utilisable. Elle rappelle aussi que la résilience ne se limite pas à une preuve mathématique : elle dépend de la qualité des capteurs, de la synchronisation temporelle, du logiciel embarqué, de la communication et de la capacité à diagnostiquer le système en fonctionnement.

\section*{Apports principaux}

Les apports de la thèse peuvent être résumés autour de quatre points. Premièrement, elle propose une approche de reconstruction d'attaques dans les systèmes CACC en modélisant les données corrompues comme des entrées inconnues. Cette approche permet d'aller au-delà de la détection et de fournir une information directement utile pour la compensation. Deuxièmement, elle développe un observateur généralisé capable de traiter des situations où les conditions classiques des observateurs à entrées inconnues ne sont pas satisfaites. Troisièmement, elle introduit une architecture hybride où l'apprentissage neuronal approxime les dynamiques résiduelles sans abandonner la structure de stabilité fournie par l'observateur. Quatrièmement, elle étend la résilience au cadre distribué grâce à un mécanisme de confiance qui adapte les poids de coopération dans un peloton.

Ces apports sont complémentaires. La reconstruction d'attaque traite la menace cyber-physique dans une chaîne CACC. L'observateur hybride traite l'incertitude de modèle et les dynamiques non linéaires. La confiance distribuée traite la fiabilité des informations coopératives. L'ensemble forme un cadre cohérent pour l'estimation résiliente dans les véhicules autonomes et connectés. Ce cadre repose sur une idée commune : exploiter les modèles, les mesures et les relations de cohérence pour produire des estimations fiables malgré des données imparfaites.

La thèse contribue également à rapprocher les méthodes de contrôle avancé des préoccupations industrielles. Les algorithmes sont formulés dans des structures compatibles avec le temps discret, les contraintes d'échantillonnage et les validations par simulation réaliste. Les perspectives d'implémentation embarquée sont explicitement abordées, ce qui facilite le passage vers des prototypes matériels. Cette orientation est importante car les solutions destinées aux véhicules autonomes doivent être non seulement performantes, mais aussi explicables, vérifiables et intégrables dans des architectures sûres.

\section*{Limites et perspectives}

Malgré les résultats obtenus, plusieurs limites restent ouvertes. Les validations expérimentales doivent être poursuivies sur des plateformes physiques, d'abord à échelle réduite puis, à terme, sur des véhicules réels. Les effets combinés de bruit capteur, délais variables, pertes de paquets, limitations de calcul et incertitudes de synchronisation devront être évalués de manière plus systématique. Le passage à l'embarqué nécessitera aussi une analyse fine de la charge de calcul, de la mémoire, de la robustesse numérique et de la tolérance aux défauts logiciels.

Une deuxième perspective concerne l'extension vers la perception coopérative. Les véhicules futurs n'échangeront pas seulement des états cinématiques, mais aussi des objets détectés, des nuages de points, des images ou des cartes locales. L'évaluation de la confiance devra donc s'étendre à des données plus riches et plus hétérogènes. Les principes développés dans cette thèse, fondés sur la cohérence physique et temporelle, pourront servir de base à ces extensions, mais de nouveaux critères seront nécessaires.

Une troisième perspective concerne l'intégration plus étroite entre estimation, contrôle et décision. Dans cette thèse, l'accent est mis sur l'estimation et sur son effet sur les fonctions de contrôle. Des travaux futurs pourraient utiliser explicitement les niveaux de confiance et les incertitudes estimées pour adapter les décisions de haut niveau : changement de mode de conduite, augmentation de la distance de sécurité, isolement d'un voisin, ou passage à une stratégie dégradée. Une telle intégration renforcerait le rôle de l'observateur comme module central de diagnostic et de décision.

Enfin, l'apprentissage dans les observateurs mérite d'être approfondi. Les réseaux neuronaux peuvent améliorer l'estimation des dynamiques complexes, mais leur utilisation dans un système critique exige des garanties supplémentaires. Les perspectives incluent l'apprentissage avec quantification d'incertitude, les architectures physiquement informées, la mise à jour sûre des paramètres et la certification de performances dans des domaines de fonctionnement bornés. Ces développements permettront de mieux concilier adaptabilité et sûreté.

\section*{Conclusion du résumé étendu}

Cette thèse propose un ensemble d'outils pour améliorer la résilience des véhicules autonomes et connectés par l'estimation d'état. Elle traite simultanément les cyberattaques sur les communications, les incertitudes non linéaires des modèles de véhicules et la fiabilité variable des données coopératives. Les méthodes développées reposent sur des observateurs à entrées inconnues, des architectures hybrides avec apprentissage et des mécanismes distribués de confiance. Les validations montrent que ces approches permettent de reconstruire des signaux d'attaque, d'améliorer l'estimation en présence de dynamiques non modélisées et de limiter l'influence de voisins compromis dans un peloton.

Au-delà des résultats spécifiques, le manuscrit défend une vision intégrée de l'estimation résiliente. Dans les véhicules connectés, la sûreté ne peut pas dépendre d'une seule source d'information supposée parfaite. Elle doit émerger de la comparaison entre mesures, modèles, communications et historiques de confiance. Les contributions de la thèse s'inscrivent dans cette direction et constituent une base pour le développement de capteurs logiciels embarqués capables de soutenir les futures fonctions de conduite coopérative, sûre et autonome.
